\subsection{Computation rules for time-dependent forms}
The goal of this section is to establish a few coordinate free rules for computing expressions involving the derivative or integral of time dependent differential forms in practice. We begin by discussing some elementary common cases and explain how to deal with more complex ones later on. 

\vspace{.25cm}

\noindent The Lie derivative $\mathcal{L}_X\alpha$ was defined as
\[
\frac{d}{dt}\bigg|_{t=0}\phi_t^\ast\alpha,
\]
where $\phi_t$ is the flow of $X$. We might also want to differentiate this expression at some other time $t=t_0$. The following proposition shows how this can be computed in terms of $\mathcal{L}_X\alpha$. In fact, we show a more general statement, where $\phi_t$ is the flow of a time dependent vector field.

\begin{proposition}
    \label{prop:lie_derivative_formula}
    Let $M$ be a manifold without boundary, $X_t$ a time dependent vector field on $M$ and $\alpha \in \Omega^r(M).$ Denote by $\phi_t$ the flow of $X_t$ and consider the open set $M_{t_0} = \{p\in M \mid \phi_t(p) \text{ exists for } t=t_0\}$. Then
    \[
        \frac{d}{dt}\bigg|_{t=t_0}(\phi_t^\ast\alpha)_p = \phi_{t_0}^\ast(\mathcal{L}_{X_{t_0}} \alpha)_p
    \]
    for all points $p \in M_{t_0}$.
\end{proposition}
\begin{proof}
    We begin by introducing some notation. Say $X_t$ is defined for time $t = t_1$. We then write $\phi_{t_1,t}$ for the solution of the following  nonautonomous ODE
    \[
    \frac{d}{dt}\phi_{t_1,t} = X_{t+t_1}(\phi_{t_1,t}),
    \]
    i.e. the flow of the reparametrised time-dependent vector field $X_{t+t_1}$. Then we have $\phi_{0,t} = \phi_t$ and 
    \[
    \phi_{t+t_0} = \phi_{t_0,t} \circ \phi_{0,t_0} 
    \]
    wherever this equation makes sense. Then we can compute:
    \begin{align*}
        \frac{d}{dt}\bigg|_{t=t_0}(\phi_t^\ast\alpha)_p &= \frac{d}{dt}\bigg|_{t=0}(\phi_{t+t_0}^\ast\alpha)_p\\ &= \frac{d}{dt}\bigg|_{t=0}(\phi_{0,t_0}^\ast\phi_{t_0,t}^\ast\alpha)_p \\
        &= \phi_{0,t_0}^\ast \frac{d}{dt}\bigg|_{t=0} (\phi_{t_0,t}^\ast\alpha)_p
    \end{align*}
    We claim that
    \[
        \frac{d}{dt}\bigg|_{t=0} (\phi_{t_0,t}^\ast\alpha)_p = (\mathcal{L}_{X_{t_0}}\alpha)_p
    \]
    To see this, we can explicitly express both sides in local coordinates and then use (\ref{eq:cartan_calc_further_rule_2}) from Lemma \ref{lem:cartan_calculus_further_rules} and (\ref{eq:cartan_calc_further_rule_3}) from Lemma \ref{lem:cartan_calculus_time_derivative_commutes_with_diff_ops} to reduce this to the case, where $\alpha$ is a $0$-form, which follows immediately. Essentially, this is again using the trick that two differential operators are equal if they satisfy the same Leibniz rule, commute with $d$ and agree on $0$-forms. 
\end{proof}

\begin{lemma}
    \label{lem:cartan_calculus_further_rules}
    Let $\alpha_t, \beta_t$ be smooth time-dependent families of differential forms on a manifold $M$ and $f \colon M \to N$ a smooth map. Then the following holds:
    \begin{enumerate}[(i)]
        \item $\frac{d}{d t}f^*\alpha_t = f^*\frac{d}{d t} \alpha_t$ \label{eq:cartan_calc_further_rule_1}
        \item $\frac{d}{dt}(\alpha_t\wedge \beta_t) = \frac{d}{dt}\alpha_t \wedge \beta_t + \alpha_t \wedge \frac{d}{dt}\beta_t$ \label{eq:cartan_calc_further_rule_2}
    \end{enumerate} 
\end{lemma}
\begin{proof}
    For (\ref{eq:cartan_calc_further_rule_1}) note that for a given point $p\in M$ the pullback along $f$ is given by a linear map
    \[
        f^\ast \colon\Lambda^r(T_{f(p)}^\ast N) \to \Lambda^r(T_{p}^\ast M),
    \] 
    so the statement follows directly from the chain rule. (\ref{eq:cartan_calc_further_rule_2}) follows from the ordinary Leibniz rule for smooth functions and can be seen by computing the derivative in local coordinates. We have in fact already seen this computation in the proof of Proposition \ref{prop:lie_derivative_properties}.
\end{proof}

This next Lemma shows, that the time derivative of time-dependent differential forms commutes with the differential operators from Theorem \ref{thm:cartan_calculus}.

\begin{lemma}
    \label{lem:cartan_calculus_time_derivative_commutes_with_diff_ops}
    Let $\alpha_t$ be a time-dependent family of differential forms and $X$ a vector field on a manifold $M$. Then the following holds:
    \begin{enumerate}[(i)]
        \item $\frac{d}{dt}d\alpha_t = d\frac{d}{dt}\alpha_t$ \label{eq:cartan_calc_further_rule_3}
        \item $\frac{d}{dt} \mathcal{L}_X \alpha_t = \mathcal{L}_X \frac{d}{dt}\alpha_t$ \label{eq:cartan_calc_further_rule_4}
        \item $\frac{d}{dt} \iota_X \alpha_t = \iota_X \frac{d}{dt}\alpha_t$ \label{eq:cartan_calc_further_rule_5}
    \end{enumerate}
\end{lemma}
\begin{proof}
    (\ref{eq:cartan_calc_further_rule_3}) follows from Schwarz's Theorem by computing the expression in local coordinates. This computation has essentially already been carried out in the proof of Proposition \ref{prop:lie_derivative_properties}. (\ref{eq:cartan_calc_further_rule_4}) follows from (\ref{eq:cartan_calc_further_rule_3}) and (\ref{eq:cartan_calc_further_rule_5}) using Cartan's magic formula. For (\ref{eq:cartan_calc_further_rule_5}) note that for $p \in M$ the map 
    \[
    \iota_{X_p} \colon \Lambda^r(T_{p}^\ast M) \to \Lambda^{r-1}(T_{p}^\ast M)
    \]
    is linear so the statement follows again from the chain rule.
\end{proof}

Moreover, the integral with respect to time also commutes with the differential operators, as the following Lemma shows.

\begin{lemma}
    Let $\alpha_t$ be a time-dependent family of differential forms defined on some interval $I$ and $X$ a vector field on a manifold $M$. Given $a,b \in I$ the following holds: 
    \begin{enumerate}[(i)]
        \item $d \int_{a}^{b} \alpha_t dt = \int_{a}^{b} d\alpha_t dt$ \label{eq:cartan_calc_further_rule_6}
        \item $\mathcal{L}_X \int_{a}^{b} \alpha_t dt = \int_{a}^{b} \mathcal{L}_X\alpha_t dt$\label{eq:cartan_calc_further_rule_7}
        \item $\iota_X \int_{a}^{b} \alpha_t dt = \int_{a}^{b} \iota_X\alpha_t dt$\label{eq:cartan_calc_further_rule_8}
    \end{enumerate}
\end{lemma}
\begin{proof}
    For (\ref{eq:cartan_calc_further_rule_6}) choose local coordinates $x_1,\ldots,x_n$ and let $\alpha^{i_1,\ldots,i_r}(t,x)$ be a component function of $\alpha_t$ with respect to these coordinates. Then 
    \[
    \frac{\partial}{\partial x_i} \int_{a}^{b} \alpha^{i_1,\ldots,i_r}(t,x) dt = \int_{a}^{b}  \frac{\partial}{\partial x_i} \alpha^{i_1,\ldots,i_r}(t,x) dt
    \]
    since $[a,b]$ is compact. The statement follows. Again (\ref{eq:cartan_calc_further_rule_7}) follows from (\ref{eq:cartan_calc_further_rule_6}) and (\ref{eq:cartan_calc_further_rule_8}) using Cartan's magic formula. (\ref{eq:cartan_calc_further_rule_8}) follows from linearity of the integral (see Lemma \ref{lem:linearity_of_integral}) by again considering the linear map $\iota_{X_p}$ for $p \in M$.
\end{proof}

Finally, we explain how to deal with time dependent vector fields in the differential operators.

\begin{lemma}
    Let $X_t$ be a time-dependent vector field on a manifold $M$. Then the following holds: 
    \begin{enumerate}[(i)]
        \item $\frac{d}{dt} \mathcal{L}_{X_t} = \mathcal{L}_{\frac{d}{dt}X_t}$ \label{eq:cartan_calc_further_rule_9}
        \item $\frac{d}{dt} \iota_{X_t} = \iota_{\frac{d}{dt}X_t}$ \label{eq:cartan_calc_further_rule_10}
    \end{enumerate}
\end{lemma}
\begin{proof}
    (\ref{eq:cartan_calc_further_rule_9}) follows from (\ref{eq:cartan_calc_further_rule_10}) as well as Lemma \ref{lem:cartan_calculus_time_derivative_commutes_with_diff_ops} (\ref{eq:cartan_calc_further_rule_3}) using Cartan's magic formula. For (\ref{eq:cartan_calc_further_rule_10}) let $\alpha \in \Omega^r(M)$ and consider the linear map 
    \[
    A \colon T_p M \to \Lambda^{r-1}(T_p^\ast M), \;\, v \mapsto \alpha_p(v,-)
    \]
    for $p \in M$. Then 
    \[
    \left( \frac{d}{dt} \iota_{X_t} \alpha \right)_p = \frac{d}{dt} A(X_t(p))
    \]
    and so the statement follows again from the chain rule.
\end{proof}

\begin{remark}
    Of course analogous statements hold for $\R^n$-valued time parameters $t\in\R^n$.
\end{remark}

\headline{Moser`s trick}

In many cases one has to differentiate a time-dependent family of differential forms, which is given by an expression where the time variable appears at multiple instances. To understand how a formula for this is obtained, we begin by discussing an important example, namely Moser`s trick, which is used to proof the Darboux Theorem about symplectic/contact manfiolds. Following that, we present a straightforward cooking recipe on how to compute the derivative of such an expression in general, illustrated by another example.

Let $M$ be a manifold (without boundary). Given a smooth family $\alpha_t$ of differential forms on $M$ indexed over $t \in \R$ we want to find a family of diffeomorphisms $\varphi_t \colon M \to M$ (i.e. an isotopy), such that $\varphi_0 = id_M$ and 
\[
    \varphi_t^\ast \alpha_t = \alpha_0 \: \text{ for all } t \in [0,1],
\]
or equivalently:
\begin{equation}
    \label{eq:mosers_trick}
    \frac{d}{dt} \varphi_t^\ast \alpha_t = 0.    
\end{equation}


% \begin{remark}
%     If $\partial M \neq \emptyset$, we call $\varphi_t$ a smooth family of diffeomorphisms if the map 
%     \[
%     M \times [0,1] \to M, \: (p,t) \mapsto \varphi_t(p)
%     \]
%     is smooth in the sense of Definition \ref{def:extension_smooth}, where we consider $M \times [0,1] \subseteq M \times \R$.
    % If $\partial M \neq \emptyset$ we say that $\varphi_t$ is a smooth family of diffeomorphisms if it can be extended to a smooth family $\varphi^\prime_t$ indexed over $\R$. In fact, if $\partial M = \emptyset$ we can also extend $\varphi_t$ to $\R$, using the following trick: The isotopy $\varphi_t$ corresponds to the level preserving diffeomorphism
    % \[
    % \Phi \colon M \times [0,1] \to M \times [0,1], (p,t) \mapsto (\varphi_t(p),t).
    % \]
    % This diffeomorphism then defines a vector field $Y:= \Phi_\ast \frac{\partial}{\partial t}$, whose integral curves uniquely determine $\varphi_t$ (this trick is also explained in \cite{Milnor1965hcobord} Thm. 5.6). We can extend this vector field arbitrarily to $M \times \R$ and then set the component pointing in the $\R$ direction to $\frac{\partial}{\partial t}$. 
%\end{remark}

Directly constructing such an isotopy is difficult but we can make use of the following trick: The isotopy $\varphi_t$ corresponds to the level preserving diffeomorphism
\[
\Phi \colon M \times [0,1] \to M \times [0,1], (p,t) \mapsto (\varphi_t(p),t).
\]
This diffeomorphism then defines a vector field $\Phi_\ast \frac{\partial}{\partial t}$ on $M \times [0,1]$, whose integral curves uniquely determine $\varphi_t$. This suggests that in order to construct an isotopy $\varphi_t$ we can construct a vector field $Y$ on $M \times [0,1]$, of the form 
\[
Y_{(p,t)} = ((X_t)_p, \frac{\partial}{\partial t}) 
\]
for a time dependent vector field $X_t$ on $M$ (this trick is also explained in \cite{Milnor1965hcobord} Thm. 5.6). 

Then (\ref{eq:mosers_trick}) can be reformulated as a condition on $X_t$, which is shown in the following Lemma.

\begin{lemma}
    \label{lem:mosers_trick_lemma}
    In the situation above the following holds:
    \[
    \frac{d}{dt} \varphi_t^\ast \alpha_t = \varphi_t^\ast(L_{X_t}\alpha_t + \dot{\alpha_t})
    \]
\end{lemma}

\begin{proof}
    The idea is to "decouple" the time variables and consider the two parameter family
    \[
        g \colon [0,1] \times \R \to T_p^\ast M, \: (t,s) \mapsto (\varphi_t^\ast \alpha_s)_p
    \]
    The partial derivatives are now much easier to compute. By Proposition \ref{prop:lie_derivative_formula} we have:
    \[
    \frac{\partial}{\partial t}\bigg|_{t=t_0} \varphi_t^\ast \alpha_s = \varphi_{t_0}^\ast(\mathcal{L}_{X_{t_0}} \alpha_s)
    \]
    Moreover, by Lemma \ref{lem:cartan_calculus_further_rules}
    \[
        \frac{\partial}{\partial s}\bigg|_{s=t_0} \varphi_t^\ast \alpha_s = \varphi_t^\ast \left(\frac{\partial}{\partial s}\bigg|_{s=t_0}  \alpha_s \right) = \varphi_t^\ast \dot{\alpha}_{t_0}.
    \]
    In order to compute the derivative of 
    \[
    f \colon [0,1] \to T_p^\ast M, \: t \mapsto (\varphi_t^\ast \alpha_t)_p
    \]
    consider the diagonal map
    \[
    \Delta \colon [0,1] \to [0,1] \times \R, \: t \mapsto (t,t).
    \]
    Then
    \[
    f = g \circ\Delta,
    \]
    hence
    \begin{align*}
        \frac{d}{dt} \bigg|_{t = t_0} f(t) & = D_{t_0}f(1) = D_{(t_0,t_0)}g \circ D_{t_0}\Delta (1) \\
        & = D_{(t_0,t_0)}(1,1) = D_{(t_0,t_0)}(1,0) + D_{(t_0,t_0)}(0,1) \\
        & = \frac{\partial g}{\partial t} (t_0,t_0) + \frac{\partial g}{\partial s} (t_0,t_0)  
    \end{align*}
    The Lemma follows.
\end{proof}
So (\ref{eq:mosers_trick}) can be reformulated to the condition:
\[
    0 = L_{X_t}\alpha_t + \dot{\alpha_t} = d\iota_{X_t}\alpha_t + \iota_{X_t}d\alpha_t + \dot{\alpha_t},
\]
where we used Theorem \ref{thm:cartan_calculus} (\ref{eq:cartan3}) in the second equation. Hence it suffices to construct a time dependent vector field $X_t$ satisfying the above conditon and whose integral curves exist for time $t = 1$. 

\headline{A cooking recipe}

The proof of Lemma \ref{lem:mosers_trick_lemma} now gives the announced cooking recipe for computing derivatives of smooth families of differential forms indexed over some time $t$, where $t$ appears at multiple instances. Namely, fix $t$ for all but one instance and differentiate. Then take the sum of all of the results. We illustrate this by an example.

\begin{example}
    Let $X_t$ be some time dependent vector field on a manifold $M$ with flow $\phi_t$. Furthermore, let $\alpha_t$ be a smooth family of differential forms and $f_t$ a smooth family of functions on $M$. We want to compute:
    \[
    \frac{d}{dt}\phi_t^\ast(e^{f_t}\alpha_t).
    \]
    Begin by fixing all but the first occurence of $t$ and differentiate, i.e. compute
    \begin{align*}
        \frac{\partial}{\partial t} \phi_t^\ast(e^{f_s}\alpha_s) = \phi_t^\ast(\mathcal{L}_{X_t}(e^{f_s}\alpha_s)) &= \phi_t^\ast((X_t \cdot e^{f_s})\alpha_s + e^{f_s}\mathcal{L}_{X_t}\alpha_s)\\ &= \phi_t^\ast((X_t \cdot f_s) e^{f_s} \alpha_s + e^{f_s}\mathcal{L}_{X_t}\alpha_s)
    \end{align*}
    Here we used Proposition \ref{prop:lie_derivative_formula} and the Leibniz rule for the Lie derivative. Next fix all but the second occurence of $t$ and using Lemma \ref{lem:cartan_calculus_further_rules} compute:
    \[
        \frac{\partial}{\partial t} \phi_s^\ast(e^{f_t}\alpha_s) = \phi_s^\ast(\frac{\partial }{\partial t}e^{f_t}\alpha_s) = \phi_s^\ast(\dot{f_t}e^{f_t}\alpha_s)
    \]
    And lastly
    \[
        \frac{\partial }{\partial t} \phi_s^\ast(e^{f_s}\alpha_t) = \phi_s^\ast(e^{f_s}\frac{\partial }{\partial t} \alpha_t) = \phi_s^\ast(e^{f_s}\dot{\alpha_t})
    \]
    Finally plug in $s=t$ and sum the results:
    \begin{align*}
        \frac{d}{dt}\phi_t^\ast(e^{f_t}\alpha_t) &= \phi_t^\ast((X_t \cdot f_t) e^{f_t} \alpha_t + e^{f_t}\mathcal{L}_{X_t}\alpha_t) + \phi_t^\ast(\dot{f_t}e^{f_t}\alpha_t) + \phi_t^\ast(e^{f_t}\dot{\alpha_t})\\
        &= \phi_t^\ast(e^{f_t}((X_t \cdot f_t) \alpha_t + \mathcal{L}_{X_t}\alpha_t + \dot{f_t}\alpha_t + \dot{\alpha_t}))
    \end{align*}
\end{example}